\chapter{用語集}\label{app:glossary}

\begin{description}\small
\item[AMBAC] Active Mass Balance Auto Control（能動的質量移動による自動姿勢制御）。手足の質量移動の反作用で姿勢を制御する考え方。『機動戦士ガンダム』の SF 考証で示された \cite{R42}。
\item[角運動量保存] 外力のモーメントがなければ、系の重心まわりの角運動量は変わらないという法則。AMBAC の出発点。
\item[機械的接続（接続）] 形状を変えたときに向きがどれだけ変わるかを与える規則 $\Gamma$。$\dot\theta=-\Gamma\cdot\dot q$。
\item[曲率] 形状空間の各点で、小さなループを回したときに向きがどれだけ進むかの密度 $F$。
\item[幾何学的位相（ホロノミー）] 形状を元に戻したのに、向きなどの量が元に戻らずに残るずれ。
\item[ホロノミック／非ホロノミック] 速度の拘束が位置の拘束に積分できる／できないこと。手足 2 本以上の AMBAC は非ホロノミック。
\item[リー括弧] 2 つの動きを交互に行ったときに生じる、どちらの方向でもない正味の動き。非ホロノミック系の可制御性の鍵。
\item[形状空間] 関節角など、機体の「形」だけを表す変数の空間。
\item[リアクションホイール（RW）] 人工衛星の姿勢制御に使う、回転させて反作用で衛星を回す円盤。AMBAC の「曲率ゼロ・連続回転できる関節」の場合にあたる。
\item[CMG] コントロール・モーメント・ジャイロ。回転するホイールの軸の向きを変えて大きなトルクを得る装置。
\item[運動量飽和] リアクションホイールが最大回転数に達し、それ以上角運動量を吸収できなくなること。AMBAC では周回周波数の上限に相当する。
\item[アンローディング] 飽和したホイールなどに溜まった角運動量を、スラスタなどの外力で逃がすこと。
\item[二乗三乗則] 大きさを $s$ 倍にすると、面積は $s^2$、体積（質量）は $s^3$ で増えるため、大型化すると相対的に力が足りなくなるという法則 \cite{R33}。
\item[フルード則] 重力が効く運動で、加速度と重力の比（フルード数）を同じに保てば、大きさが違っても同じように振る舞うという相似則 \cite{R29}。
\item[ZMP] Zero Moment Point。地面の反力の合力が通る点。接地範囲の外に出ようとすると転倒が始まる \cite{R28}。
\item[支持範囲（支持多角形）] 接地している部分を囲む範囲。ZMP はこの中にしか来られない。
\item[スティックスリップ] 静止摩擦で止まる状態と、動摩擦で滑る状態を繰り返す現象 \cite{R30}。
\item[腰の戦略（ヒップ戦略）] 上体を前後に動かしてバランスを取る姿勢制御の方法 \cite{R31}。足首の戦略（足首の回転でバランスを取る）と対になる。
\item[球面座標] 方向を方位角 $\psi$ と仰角 $\alpha$ で表す座標。2 軸の肩の関節角にあたる。
\item[ジンバルロック] 2 つ以上の回転軸が揃って自由度を失う現象。球面座標では極（真上・真下）で起き、関節速度が発散する。
\item[コッドマンのパラドックス] 腕をひねらずに動かして元の向きに戻すと、上腕が長軸まわりにひねられている現象 \cite{R34}。ひねり＝腕の向きが囲んだ立体角。
\item[立体角] 球面上の領域の面積（単位球で測ったもの）。全球は $4\pi$。
\item[ガウス・ボンネの定理] 曲面上で閉曲線に沿って平行移動したベクトルの回転角が、囲まれた領域の曲率の積分に等しいという定理 \cite{R11}。
\item[上肢帯] 鎖骨と肩甲骨。腕の付け根（肩関節）そのものを動かす。
\item[前方突出（protraction）] 肩甲骨・鎖骨の動きで肩を前に出すこと。
\item[内旋・外旋] 腕を長軸まわりに内側・外側へひねること。
\item[水平内転] 横に上げた腕を、水平面内で体の前へ寄せること。
\item[線形計画（LP）] 目的関数と制約がすべて 1 次式の最適化問題。確実に大域解が求まる。
\item[SQP（SLSQP）] 逐次二次計画法。非線形の制約付き最適化の解法 \cite{R26}。
\item[モデル予測制御（MPC）] 一定時間ごとに最適制御問題を解き直し、最初の部分だけを実行する制御。
\item[進化戦略（ES）] 方策のパラメータにノイズを加えて試し、報酬の良い方向へ動かす最適化法 \cite{R19}。
\item[帰納バイアス] 学習に与える前提知識・構造。方策の形などとして与える。
\item[CPG] Central Pattern Generator（中枢パターン生成器）。リズミカルな運動を生む仕組み \cite{R22}。
\item[PMTG] Policies Modulating Trajectory Generators。軌道生成器を学習した方策で調整する構造 \cite{R21}。
\item[ドメインランダム化] 学習時にシミュレーションのパラメータを乱数で変え、モデルの誤差に強い方策を得る方法 \cite{R23}。
\item[最小躍度軌道] 躍度（加速度の時間変化）の 2 乗の積分を最小にする滑らかな軌道。始点と終点で速度・加速度がゼロ \cite{R27}。
\item[準静的] 慣性力が無視できるほどゆっくり動くとみなす扱い。
\item[アセンブル] パーツを組み合わせて機体を構成すること（\cite{R3} にならった呼び方）。第 2 段階の枠組みの考え方。
\item[校正値] シミュレーションと実機の差を埋めるために測定・推定して与える値（慣性、摩擦、可動域、アクチュエータ特性など）。
\end{description}
